\documentclass[a4paper,12pt]{article}

% Быстрый вариант для Overleaf: использовать XeLaTeX.
% В Overleaf: Menu -> Compiler -> XeLaTeX.
% Это убирает медленную генерацию кириллических шрифтов pdfLaTeX.
\usepackage{fontspec}
\usepackage[russian]{babel}
\setmainfont{Liberation Serif}
\setsansfont{Liberation Sans}
\setmonofont{Liberation Mono}

\usepackage{amsmath,amssymb}
\usepackage{graphicx}
\usepackage{float}
\usepackage{booktabs}
\usepackage{array}
\usepackage{geometry}
\usepackage[hidelinks]{hyperref}
\usepackage{caption}
\usepackage{placeins}
\usepackage{indentfirst}

\geometry{left=25mm,right=15mm,top=20mm,bottom=20mm}
\setlength{\parindent}{1.25cm}
\setlength{\parskip}{0.2em}

% Папки с рисунками. Основной вариант — figure_lab3.
% Если папка называется figures_lab3, файл тоже будет найден.
\newcommand{\missingfigurebox}[1]{%
  \fbox{\parbox{0.86\linewidth}{\centering Файл рисунка \texttt{\detokenize{#1.png}} не найден.\\
  Проверьте название файла и папки с рисунками.}}%
}

\newcommand{\includeLabGraphic}[2][]{%
  \IfFileExists{figure_lab3/#2.png}{\includegraphics[#1]{figure_lab3/#2.png}}{%
    \IfFileExists{figures_lab3/#2.png}{\includegraphics[#1]{figures_lab3/#2.png}}{%
      \missingfigurebox{#2}%
    }%
  }%
}

\newcommand{\labfigure}[4][0.92\linewidth]{%
\begin{figure}[htbp]
    \centering
    \includeLabGraphic[width=#1]{#2}
    \caption{#3}
    \label{#4}
\end{figure}
}

\newcommand{\optionallabfigure}[4][0.92\linewidth]{%
\IfFileExists{figure_lab3/#2.png}{%
\begin{figure}[htbp]
    \centering
    \includegraphics[width=#1]{figure_lab3/#2.png}
    \caption{#3}
    \label{#4}
\end{figure}}{%
\IfFileExists{figures_lab3/#2.png}{%
\begin{figure}[htbp]
    \centering
    \includegraphics[width=#1]{figures_lab3/#2.png}
    \caption{#3}
    \label{#4}
\end{figure}}{}%
}%
}

\begin{document}

\begin{titlepage}
\begin{center}
    \large
    Федеральное государственное автономное образовательное учреждение высшего образования\\
    \textbf{Московский физико-тезнический институт}\\[0.3cm]
    Физтех-школа природоподобных, плазменных и ядерных технологий\\
    имени И. В. Курчатова

    \vfill

    \textbf{Отчёт по лабораторной работе №7}\\[0.4cm]
    \Large\textbf{Обработка цифровых сигналов}\\[0.4cm]
    \normalsize
    Дискретное преобразование Фурье, вейвлет-анализ, КИХ- и БИХ-фильтрация, оценка ЧСС

    \vfill

    \begin{flushright}
        \textbf{Выполнил:} студент группы Б07-202 \\[0.2cm]
        \textbf{ФИО:Андрулис Валерий Виталевич} \\[0.2cm]
        \textbf{Вариант:} 3\\[0.2cm]
        \textbf{Преподаватель: Карпов Валерий Эдуардович}
    \end{flushright}

    \vfill

    \noindent\textbf{GitHub-репозиторий с исходными кодами:}\\
    \href{https://github.com/Venzoll-research/Robolab7_signals}{\texttt{Venzoll-research/Robolab7\_signals}}

    \vfill

    Москва, 2026
\end{center}
\end{titlepage}

\tableofcontents
\newpage

\section{Цель и постановка работы}

Цель лабораторной работы состоит в освоении базовых математических моделей и алгоритмов цифровой обработки сигналов: дискретного преобразования Фурье, частотно-временного анализа и цифровой фильтрации.

В работе использовались два набора данных. Первый набор --- дискретный сигнал варианта 3 из файла \texttt{signal\_3.txt}. Для него частота дискретизации была принята равной
\[
    f_s = 40~\text{Гц}, \qquad T = \frac{1}{f_s}=0.025~\text{с}.
\]
Количество отсчётов равно $N=400$, длительность записи составляет $10$~с. Второй набор данных --- зашумлённый ЭКГ-сигнал \texttt{signal\_ecg\_3.txt}; для него использована частота дискретизации $f_s=360$~Гц.

Согласно варианту 3 для фильтрации основного сигнала должны быть построены режекторные фильтры: КИХ-фильтр с окном Блэкмана и БИХ-фильтр Чебышёва I рода. Режекторный фильтр подавляет выбранную полосу частот и сохраняет частотные компоненты вне этой полосы.

\section{Краткие теоретические сведения}

\subsection{Представления цифрового сигнала}

Сигналом называется физическая величина, изменяющаяся во времени и несущая информацию о состоянии системы. Временное представление показывает форму сигнала, но не всегда позволяет понять, какие частотные компоненты в нём присутствуют. Поэтому в цифровой обработке применяются дополнительные представления: частотное, комплексное, $z$-представление и частотно-временное представление.

Частотное представление используется для спектрального анализа, выделения гармоник и проектирования фильтров. Вейвлет-представление позволяет анализировать не только частоту, но и момент появления частотной компоненты, что особенно важно для нестационарных сигналов.

\subsection{Дискретное преобразование Фурье}

Для последовательности $x[n]$, $n=0,1,\ldots,N-1$, дискретное преобразование Фурье определяется выражением
\begin{equation}
    X[m] = \sum_{n=0}^{N-1} x[n] e^{-j\frac{2\pi nm}{N}}, \qquad m=0,1,\ldots,N-1.
    \label{eq:dft}
\end{equation}

Эквивалентно его можно записать через действительную и мнимую части:
\begin{equation}
    X[m] = \sum_{n=0}^{N-1} x[n]\left(\cos\frac{2\pi nm}{N} - j\sin\frac{2\pi nm}{N}\right).
\end{equation}

Обратное ДПФ восстанавливает временную последовательность по спектру:
\begin{equation}
    x[n] = \frac{1}{N}\sum_{m=0}^{N-1} X[m] e^{j\frac{2\pi nm}{N}}.
    \label{eq:idft}
\end{equation}

Для вещественного сигнала спектр симметричен: модули положительных и отрицательных частот совпадают, а фазы имеют противоположные знаки. Поэтому для анализа достаточно использовать односторонний спектр до частоты Найквиста
\begin{equation}
    f_N = \frac{f_s}{2}.
\end{equation}

При построении одностороннего амплитудного спектра постоянная составляющая не удваивается, а амплитуды внутренних положительных частот умножаются на два:
\begin{equation}
    A[k] = \frac{2}{N}|X[k]|, \qquad k=1,2,\ldots,\frac{N}{2}-1.
\end{equation}

Для проверки корректности реализации ДПФ используется сравнение с библиотечным БПФ. В выполненной программе максимальное расхождение между ручным ДПФ и \texttt{numpy.fft} составило
\[
    \max |X_{\text{manual}} - X_{\text{fft}}| \approx 3.50\cdot 10^{-10},
\]
что соответствует ошибке численного округления.

\subsection{Вейвлет-преобразование}

Если частотный состав сигнала меняется во времени, обычный спектр Фурье уже не показывает, когда именно появляется та или иная гармоника. В этом случае применяется вейвлет-преобразование. Непрерывное вейвлет-преобразование имеет вид
\begin{equation}
    W(a,b) = \frac{1}{\sqrt{a}}\int_{-\infty}^{+\infty} x(t)\psi^*\left(\frac{t-b}{a}\right)dt,
    \label{eq:cwt}
\end{equation}
где $a$ --- масштаб, $b$ --- сдвиг, $\psi(t)$ --- материнский вейвлет. Чем меньше масштаб, тем выше анализируемая частота; чем больше масштаб, тем ниже частота.

В работе использовалось непрерывное вейвлет-преобразование CWT с комплексным вейвлетом \texttt{cmor1.5-1.0}. Скалограмма строилась как модуль коэффициентов CWT.

\subsection{Линейная цифровая фильтрация}

Линейный цифровой фильтр в общем случае описывается разностным уравнением
\begin{equation}
    y[n] = \frac{1}{a_0}\left(\sum_{i=0}^{M-1} b_i x[n-i] - \sum_{i=1}^{P-1} a_i y[n-i]\right),
    \label{eq:difference}
\end{equation}
где $b_i$ --- коэффициенты прямой связи, $a_i$ --- коэффициенты обратной связи. Если обратной связи нет, фильтр является КИХ-фильтром. Если обратная связь присутствует, фильтр является БИХ-фильтром.

КИХ-фильтры устойчивы по конструкции, так как их импульсная характеристика конечна. БИХ-фильтры обычно требуют меньшего числа коэффициентов для достижения резкого изменения АЧХ, но их устойчивость зависит от расположения полюсов в $z$-плоскости.

\subsection{КИХ-фильтр и окно Блэкмана}

КИХ-фильтр можно получить оконным методом. Сначала задаётся идеальная импульсная характеристика фильтра, затем она умножается на конечное окно. Окно Блэкмана задаётся формулой
\begin{equation}
    w[n] = 0.42 - 0.5\cos\left(\frac{2\pi n}{N-1}\right) + 0.08\cos\left(\frac{4\pi n}{N-1}\right),
    \qquad n=0,1,\ldots,N-1.
    \label{eq:blackman}
\end{equation}

Окно уменьшает боковые лепестки частотной характеристики, но расширяет переходную полосу. В работе был построен КИХ-режектор с окном Блэкмана, числом коэффициентов $61$, центральной частотой режекции $5$~Гц и шириной полосы $0.8$~Гц.

\subsection{БИХ-фильтр Чебышёва I рода}

Фильтр Чебышёва I рода имеет пульсации в полосе пропускания и более резкий переход между полосами по сравнению с фильтром Баттерворта того же порядка. В работе использован режекторный БИХ-фильтр Чебышёва I рода первого порядка с пульсацией $0.5$~дБ и полосой подавления $4.5$--$5.5$~Гц.

Передаточная функция цифрового фильтра может быть записана как
\begin{equation}
    H(z)=\frac{b_0+b_1z^{-1}+\ldots+b_Mz^{-M}}{a_0+a_1z^{-1}+\ldots+a_Pz^{-P}}.
\end{equation}

Для анализа устойчивости БИХ-фильтра строится карта нулей и полюсов. Цифровой БИХ-фильтр устойчив, если все полюса лежат внутри единичной окружности на комплексной $z$-плоскости.

\subsection{Оценка частоты сердечных сокращений}

Для зашумлённого ЭКГ-сигнала требуется определить среднюю частоту сердечных сокращений. Сначала сигнал фильтруется в полосе $5$--$15$~Гц, где наиболее выражены QRS-комплексы. Затем находятся R-пики. Средняя частота сердечных сокращений рассчитывается по среднему RR-интервалу:
\begin{equation}
    \text{ЧСС} = \frac{60 f_s}{RR_{\text{отсчёты}}}.
    \label{eq:hr}
\end{equation}

\section{Ручной расчёт ДПФ и спектральный анализ}

В программе ДПФ реализовано вручную через матрицу комплексных экспонент. Для каждого индекса $m$ вычислялась сумма \eqref{eq:dft}. Библиотечное БПФ использовалось только для проверки правильности результата.

На рис.~\ref{fig:dft} приведены амплитудный спектр, действительная часть, мнимая часть и фазовый спектр. При построении фазы и комплексных частей малые спектральные значения были отброшены маской по амплитуде, так как фаза около нулевых амплитуд не имеет физического смысла и определяется численным шумом.

\labfigure{manual_dft_spectrum}{Результат ручного ДПФ: амплитудный спектр, действительная часть, мнимая часть и фаза.}{fig:dft}

По амплитудному спектру были выделены основные гармоники. Полученное аналитическое представление сигнала в синусной форме:
\begin{align}
    x(t) \approx{}& 3.99\sin(2\pi\cdot 1.00t)
    + 12.44\sin(2\pi\cdot 2.00t)
    + 17.90\sin(2\pi\cdot 3.00t) \notag\\
    &+ 1.80\sin(2\pi\cdot 6.00t)
    + 3.35\sin(2\pi\cdot 9.00t)
    + 1.71\sin(2\pi\cdot 10.00t) \notag\\
    &+ 15.28\sin(2\pi\cdot 14.00t).
    \label{eq:signal-model}
\end{align}

Восстановленный по найденным гармоникам сигнал практически совпадает с исходным. Это подтверждает, что основные гармоники были определены корректно.

\labfigure{signal_reconstruction}{Сравнение исходного сигнала и сигнала, восстановленного по аналитическому представлению.}{fig:reconstruction}

\optionallabfigure{signal_reconstruction_error}{Ошибка восстановления сигнала по найденным гармоникам.}{fig:reconstruction-error}

\FloatBarrier

\section{Нестационарный сигнал и вейвлет-анализ}

Для проверки частотно-временного анализа был сформирован модифицированный сигнал. В отличие от исходного стационарного сигнала, его гармоники существуют не на всей длительности записи, а только на отдельных временных интервалах. Это позволяет показать ограничение обычного спектра Фурье: он хорошо определяет набор частот, но не показывает время их появления.

\labfigure{stationary_and_local_signal}{Сравнение исходного стационарного сигнала и сигнала с локальными синусоидами.}{fig:stationary-local}

Далее для исходного и модифицированного сигналов были построены скалограммы. По скалограмме стационарного сигнала частотные компоненты наблюдаются на всём интервале времени. В модифицированном сигнале частотные полосы становятся локальными, что соответствует заданным временным окнам активности гармоник.

\labfigure{wavelet_scalograms}{Скалограммы исходного и модифицированного сигналов.}{fig:wavelets}

\FloatBarrier

\section{Проектирование и применение режекторных фильтров}

\subsection{КИХ-режектор с окном Блэкмана}

КИХ-фильтр был построен оконным методом. Идеальная импульсная характеристика режекторного фильтра была ограничена окном Блэкмана \eqref{eq:blackman}. Для КИХ-фильтра коэффициент обратной связи отсутствует, поэтому в программе использовалось
\[
    a_{\text{FIR}} = [1].
\]
Фильтрация выполнялась вручную по разностному уравнению \eqref{eq:difference}; для КИХ-фильтра это уравнение сводится к свёртке входного сигнала с коэффициентами $b_i$.

\optionallabfigure[0.75\linewidth]{fir_blackman_window}{Окно Блэкмана, использованное при проектировании КИХ-фильтра.}{fig:blackman-window}

\labfigure{fir_blackman_response}{АЧХ и ФЧХ КИХ-режектора с окном Блэкмана.}{fig:fir-response}

\subsection{БИХ-режектор Чебышёва I рода}

Для сравнения был построен БИХ-фильтр Чебышёва I рода. Коэффициенты фильтра получены функцией \texttt{cheby1}, а сама фильтрация выполнена вручную по разностному уравнению. Для используемых параметров были получены коэффициенты
\begin{equation}
    b = [0.97324415,\; -1.38063110,\; 0.97324415],
\end{equation}
\begin{equation}
    a = [1.00000000,\; -1.38063110,\; 0.94648830].
\end{equation}

\labfigure{iir_cheby1_response}{АЧХ и ФЧХ БИХ-режектора Чебышёва I рода.}{fig:iir-response}

На карте нулей и полюсов нули располагаются около единичной окружности в области подавляемой частоты, формируя провал АЧХ. Полюса находятся внутри единичной окружности, что соответствует устойчивости цифрового БИХ-фильтра.

\labfigure[0.72\linewidth]{iir_cheby1_zero_pole_map}{Карта нулей и полюсов БИХ-режектора.}{fig:iir-zpk}

\subsection{Сравнение КИХ- и БИХ-фильтрации}

Оба фильтра были применены к исходному сигналу. Результаты показаны на рис.~\ref{fig:filter-comparison}. Временные реализации после фильтрации близки к исходному сигналу, так как выбранная полоса подавления $4.5$--$5.5$~Гц не содержит самой мощной гармоники исходного сигнала. При этом фильтры демонстрируют различное вычислительное время.

\labfigure{filtering_comparison_time}{Сравнение результатов КИХ- и БИХ-фильтрации во времени.}{fig:filter-comparison}

В одном из запусков были получены следующие времена ручной фильтрации:
\begin{table}[H]
\centering
\caption{Сравнение времени ручной фильтрации}
\label{tab:filter-time}
\begin{tabular}{lcc}
\toprule
Фильтр & Время, мкс & Особенность реализации \\
\midrule
КИХ, окно Блэкмана & $3369.3$ & 61 коэффициент, нет обратной связи \\
БИХ, Чебышёв I рода & $412.7$ & 3 коэффициента числителя и 3 коэффициента знаменателя \\
\bottomrule
\end{tabular}
\end{table}

БИХ-фильтр оказался быстрее, так как для его вычисления требуется существенно меньше коэффициентов. Это соответствует общему свойству БИХ-фильтров: они часто достигают требуемой частотной характеристики меньшим порядком, чем КИХ-фильтры.

\optionallabfigure{fir_filtered_signal}{Результат КИХ-фильтрации исходного сигнала.}{fig:fir-filtered-signal}
\optionallabfigure{fir_spectrum_before_after}{Амплитудные спектры до и после КИХ-фильтрации.}{fig:fir-spectrum-before-after}
\optionallabfigure{iir_filtered_signal}{Результат БИХ-фильтрации исходного сигнала.}{fig:iir-filtered-signal}
\optionallabfigure{iir_spectrum_before_after}{Амплитудные спектры до и после БИХ-фильтрации.}{fig:iir-spectrum-before-after}

\FloatBarrier

\section{Оценка средней частоты сердечных сокращений}

В практической части был обработан зашумлённый ЭКГ-сигнал. Для подавления шума использован полосовой фильтр Баттерворта с полосой пропускания $5$--$15$~Гц. Такая полоса применяется для выделения QRS-комплексов, так как они имеют резкие фронты и хорошо выражены в данном диапазоне.

\labfigure{ecg_bandpass_response}{АЧХ полосового фильтра Баттерворта 5--15 Гц для обработки ЭКГ.}{fig:ecg-bandpass}

После фильтрации сигнал стал более пригодным для поиска R-пиков: низкочастотные и высокочастотные помехи были частично подавлены.

\labfigure{ecg_filtered_signal}{Исходный и отфильтрованный ЭКГ-сигнал.}{fig:ecg-filtered}

В текущей реализации R-пики определялись функцией \texttt{find\_peaks} на отфильтрованном сигнале. Использовались два основных ограничения: порог по высоте и минимальное расстояние между соседними пиками. Порог был задан как
\[
    h_{\text{thr}} = 0.5\max|x_{\text{filtered}}[n]|,
\]
а минимальное расстояние между пиками --- $0.3$~с, то есть
\[
    d_{\min} = 0.3 f_s.
\]
При $f_s=360$~Гц это соответствует $108$ отсчётам.

\optionallabfigure{ecg_peak_detection_steps}{Промежуточные этапы выделения R-пиков ЭКГ.}{fig:ecg-steps}

\labfigure{ecg_r_peaks}{Обнаруженные R-пики на отфильтрованном ЭКГ-сигнале.}{fig:ecg-peaks}

В результате было найдено $19$ R-пиков. Средний RR-интервал составил
\[
    RR_{\text{mean}} = 185.11~\text{отсчётов} \approx 514.2~\text{мс}.
\]
По формуле \eqref{eq:hr} получена средняя частота сердечных сокращений
\[
    \text{ЧСС} = \frac{60\cdot 360}{185.11} \approx 116.7~\text{уд/мин}.
\]

\FloatBarrier

\section{Ответы на контрольные вопросы}

\subsection*{1. Как вручную вычислить ДПФ для последовательности \texorpdfstring{$\{1,0,-1,0\}$}{\{1,0,-1,0\}}?}

Для $N=4$ используется формула \eqref{eq:dft}. Последовательность $x=[1,0,-1,0]$. Тогда
\[
X[0]=1+0-1+0=0,
\]
\[
X[1]=1+0\cdot e^{-j\pi/2}+(-1)e^{-j\pi}+0\cdot e^{-j3\pi/2}=1+1=2,
\]
\[
X[2]=1+0\cdot e^{-j\pi}+(-1)e^{-j2\pi}+0\cdot e^{-j3\pi}=1-1=0,
\]
\[
X[3]=1+0\cdot e^{-j3\pi/2}+(-1)e^{-j3\pi}+0\cdot e^{-j9\pi/2}=1+1=2.
\]
Итого $X=[0,2,0,2]$.

\subsection*{2. Почему в ОДПФ нужна нормировка?}

Прямое ДПФ суммирует вклад всех $N$ отсчётов. Чтобы обратное преобразование вернуло исходные амплитуды, в ОДПФ вводится множитель $1/N$. Без этой нормировки восстановленный сигнал был бы увеличен в $N$ раз.

\subsection*{3. Что такое КИХ- и БИХ-фильтры?}

КИХ-фильтр имеет конечную импульсную характеристику и не использует обратную связь. Его выход зависит только от текущего и предыдущих входных отсчётов. БИХ-фильтр имеет бесконечную импульсную характеристику, так как использует обратную связь: выход зависит не только от входных отсчётов, но и от предыдущих выходных значений. КИХ-фильтры устойчивы по конструкции и часто имеют линейную фазу; БИХ-фильтры обычно требуют меньшего порядка, но требуют проверки устойчивости.

\subsection*{4. Как записать разностное уравнение по передаточной функции БИХ-фильтра?}

Если
\[
    H(z)=\frac{b_0+b_1z^{-1}+\ldots+b_Mz^{-M}}{a_0+a_1z^{-1}+\ldots+a_Pz^{-P}},
\]
то соответствующее разностное уравнение имеет вид
\[
    y[n]=\frac{1}{a_0}\left(\sum_{i=0}^{M} b_i x[n-i]-\sum_{i=1}^{P} a_i y[n-i]\right).
\]
Именно такая форма была использована в ручной функции фильтрации.

\section{Выводы}

В ходе работы была реализована ручная версия дискретного преобразования Фурье. Сравнение с \texttt{numpy.fft} показало расхождение порядка $10^{-10}$, что подтверждает корректность реализации. По амплитудному спектру были выделены основные гармоники исходного сигнала и построено аналитическое представление в виде суммы синусоид. Восстановленный сигнал совпал с исходным с пренебрежимо малой ошибкой.

Для демонстрации ограничения обычного спектра Фурье был построен сигнал с локальными синусоидами. Вейвлет-скалограмма показала, что частотно-временной анализ позволяет определить не только частоты, но и интервалы времени, в течение которых эти частоты присутствуют.

Для варианта 3 были построены два режекторных фильтра: КИХ-фильтр с окном Блэкмана и БИХ-фильтр Чебышёва I рода. Фильтрация в обоих случаях выполнялась вручную по разностному уравнению. БИХ-фильтр показал меньшую вычислительную задержку благодаря меньшему числу коэффициентов. Карта нулей и полюсов подтвердила устойчивость БИХ-фильтра.

В практической части был обработан зашумлённый ЭКГ-сигнал. После полосовой фильтрации 5--15 Гц были найдены R-пики и рассчитана средняя частота сердечных сокращений. Полученное значение составило примерно $116.7$~уд/мин.

\begin{thebibliography}{9}
\bibitem{labmanual}
Леушина В. В., Сорокоумов П. С. \href{https://drive.google.com/file/d/1_ynRAScdyAZE_pU9GG6RBBCch2E7XNu5/view?usp=drive_link}{\textit{Методические материалы по выполнению практических работ по курсу «Автоматическое и интеллектуальное управление техническими системами». Работа 3. Обработка цифровых сигналов}}. 2026.
\end{thebibliography}

\end{document}
